\documentclass[11pt]{article}
\usepackage[margin=1in]{geometry}
\usepackage{amsmath,amssymb,amsthm,mathtools}
\usepackage{microtype}
\usepackage[hidelinks]{hyperref}

\newtheorem{theorem}{Theorem}
\newtheorem{lemma}{Lemma}
\newtheorem{corollary}{Corollary}
\theoremstyle{remark}
\newtheorem{remark}{Remark}

\newcommand{\T}{\mathbb T}
\newcommand{\C}{\mathbb C}
\newcommand{\A}{\mathcal A}
\newcommand{\D}{\mathcal D}
\newcommand{\wh}{\widehat}

\title{\textbf{Fractional heat generator domains on the circle are algebras}\\
\large A partial answer to Remark 2.9 of arXiv:1901.07477}
\author{Automated mathematics run \texttt{fa\_banach\_001}, lane 7}
\date{9 August 2026}

\begin{document}
\maketitle

\begin{center}
\fbox{\parbox{0.92\textwidth}{
\textbf{Status.} Candidate partial result; current verdict: \emph{likely
valid}. Human review is recommended. For every $0<\alpha\le2$, the packet
settles the source's generator-core algebra question affirmatively for the
wrapped fractional heat convolution semigroup on the classical compact group
$\T$. In fact, the entire generator domain in $C(\T)$ is an algebra. The
general locally compact quantum-group question remains open.
}}
\end{center}

\section{Source question}

Skalski and Viselter, \emph{Generating functionals for locally compact
quantum groups}, arXiv:1901.07477, Remark 2.9 on page 9, associate to a
symmetric convolution semigroup of states a completely positive contraction
semigroup $(T_t)_{t\ge0}$ on the universal quantum-group $C^*$-algebra, with
generator $L$. Their twisted Fourier cone is denoted $\D_+$. They prove that
\[
                         \operatorname{span}(\D_+)\cap D(L)
\]
is a dense selfadjoint invariant core for $L$, and then state: ``We do not
know however if it is an algebra'' \cite[p.~9, Remark~2.9]{SV}.

We answer this question for a natural one-parameter family on the circle.
Identify $\T=\mathbb R/(2\pi\mathbb Z)$. For $0<\alpha\le2$, let
$(\mu_t^{(\alpha)})_{t\ge0}$ be the symmetric convolution semigroup determined
by
\[
       \wh{\mu_t^{(\alpha)}}(n)=e^{-t|n|^\alpha},\qquad n\in\mathbb Z.
\]
Its convolution operators on $C(\T)$ are denoted $T_t^{(\alpha)}$, and the
positive generator is
\[
        A_\alpha=-L_\alpha=|D|^\alpha,\qquad
        \wh{A_\alpha f}(n)=|n|^\alpha\wh f(n).
\]

\section{Idea of the proof}

If $f\in D(A_\alpha)$ and $h=A_\alpha f$, then, modulo its constant Fourier
coefficient, $f$ is the periodic Riesz potential $K_\alpha*h$, where
$\wh K_\alpha(n)=|n|^{-\alpha}$ for $n\ne0$. Translation estimates for this
kernel give exactly the modulus of continuity
\[
 \omega_f(r)\lesssim\|h\|_\infty
 \begin{cases}
 r^\alpha,&0<\alpha<1,\\
 r\log(e/r),&\alpha=1,\\
 r,&1<\alpha<2.
 \end{cases}
\tag{1}\label{eq:modulus-intro}
\]

For $0<\alpha<2$, the product rule for the fractional Laplacian contains the
carr\'e-du-champ term
\[
 \Gamma_\alpha(f,g)(x)
 =c_\alpha\int_{-\pi}^{\pi}
   (f(x)-f(x+y))(g(x)-g(x+y))q_\alpha(y)\,dy,
\]
where $q_\alpha(y)\asymp |y|^{-1-\alpha}$ near zero.  The product of the two
moduli in \eqref{eq:modulus-intro} makes this integral absolutely and
uniformly convergent: the local majorant is respectively
$r^{\alpha-1}$, $\log^2(e/r)$, or $r^{1-\alpha}$. All are integrable in the
stated ranges. Thus $\Gamma_\alpha(f,g)$ is continuous. Poisson smoothing and
closedness of $A_\alpha$ then upgrade the formal product rule to the full
generator domain.

\section{Main result}

\begin{theorem}[Domain-algebra theorem]
\label{thm:main}
For every $0<\alpha\le2$, the domain $D(A_\alpha)$ of the generator of
$(T_t^{(\alpha)})_{t\ge0}$ on $C(\T)$ is closed under pointwise
multiplication. More precisely, if $0<\alpha<2$ and $f,g\in D(A_\alpha)$,
then
\[
 A_\alpha(fg)
 =fA_\alpha g+gA_\alpha f-\Gamma_\alpha(f,g),
\tag{2}\label{eq:product}
\]
where $\Gamma_\alpha(f,g)\in C(\T)$ is the absolutely convergent integral
above.
\end{theorem}

\begin{corollary}[Partial answer to the source]
\label{cor:source}
For the classical locally compact quantum group $\T$ and every wrapped
fractional heat semigroup above,
\[
               \operatorname{span}(\D_+)\cap D(L_\alpha)
\]
is an algebra.
\end{corollary}

\section{Proof}

We normalize convolution using Haar probability measure. Constants below may
depend on $\alpha$ but not on the functions.

\begin{lemma}[Periodic potential modulus]
\label{lem:kernel}
Let $0<\alpha<2$, and let $K_\alpha$ be the periodic distribution with
Fourier coefficients
\[
 \wh K_\alpha(0)=0,
 \qquad \wh K_\alpha(n)=|n|^{-\alpha}\quad(n\ne0).
\]
Then $K_\alpha\in L^1(\T)$ and, for $0<|y|\le1/2$,
\[
 \|K_\alpha(\cdot+y)-K_\alpha\|_1
 \le C_\alpha\rho_\alpha(|y|),
 \qquad
 \rho_\alpha(r)=
 \begin{cases}
 r^\alpha,&0<\alpha<1,\\
 r\log(e/r),&\alpha=1,\\
 r,&1<\alpha<2.
 \end{cases}
\tag{3}\label{eq:kernelmod}
\]
Consequently, if $f\in D(A_\alpha)$, then
\[
 \omega_f(r):=\sup_{|y|\le r}\|f(\cdot+y)-f\|_\infty
 \le C_\alpha\|A_\alpha f\|_\infty\rho_\alpha(r).
\tag{4}\label{eq:fmod}
\]
\end{lemma}

\begin{proof}
The standard periodic Riesz-kernel estimates, obtained for example by Abel
regularization and summation by parts, are
\[
 \begin{array}{ll}
 |K_\alpha(x)|\le C|x|^{\alpha-1},\quad
 |K_\alpha'(x)|\le C|x|^{\alpha-2},&0<\alpha<1,\\[2mm]
 K_1(x)=-2\log(2|\sin(x/2)|),\quad
 |K_1'(x)|\le C|x|^{-1},&\alpha=1,\\[2mm]
 K_\alpha'\in L^1(\T),&1<\alpha<2.
 \end{array}
\tag{5}\label{eq:kernelpoint}
\]
For $\alpha<1$, split the translation integral into
$\{|x|\le2|y|\}$ and its complement. The first part is $O(|y|^\alpha)$
by the bound on $K_\alpha$; the mean-value theorem and the derivative bound
give the same estimate on the complement. The logarithmic line of
\eqref{eq:kernelpoint} gives $O(|y|\log(e/|y|))$ in exactly the same split.
For $\alpha>1$, the $W^{1,1}$ translation estimate gives $O(|y|)$. This
proves \eqref{eq:kernelmod} and also $K_\alpha\in L^1$.

Now put $h=A_\alpha f$. Fourier coefficients, first in the distributional
sense and then in $C(\T)$, give
\[
              f-\wh f(0)=K_\alpha*h.
\]
Taking a translation difference and using Young's inequality together with
\eqref{eq:kernelmod} proves \eqref{eq:fmod}.
\end{proof}

For $0<\alpha<2$, let
\[
 q_\alpha(y)=\sum_{k\in\mathbb Z}|y+2\pi k|^{-1-\alpha},
 \qquad -\pi<y<\pi, y\ne0.
\]
Up to a positive normalization constant, the usual periodic singular-integral
formula reads
\[
 A_\alpha u(x)=c_\alpha\,\operatorname{pv}
 \int_{-\pi}^{\pi}(u(x)-u(x+y))q_\alpha(y)\,dy
\tag{6}\label{eq:singint}
\]
for smooth $u$; it is the periodization of the fractional-Laplacian formula
\cite{Stinga}. Notice that $q_\alpha(y)\le C_\alpha|y|^{-1-\alpha}$ near
zero and is bounded on compact subintervals avoiding zero.

Let $f,g\in D(A_\alpha)$. Lemma~\ref{lem:kernel} gives
\[
 \begin{aligned}
 &|(f(x)-f(x+y))(g(x)-g(x+y))q_\alpha(y)|\\
 &\qquad\le C\|A_\alpha f\|_\infty\|A_\alpha g\|_\infty
 \begin{cases}
 |y|^{\alpha-1},&0<\alpha<1,\\
 \log^2(e/|y|),&\alpha=1,\\
 |y|^{1-\alpha},&1<\alpha<2.
 \end{cases}
\end{aligned}
\tag{7}\label{eq:majorant}
\]
Each right-hand side is integrable at zero. The integral defining
$\Gamma_\alpha(f,g)$ is therefore absolutely and uniformly convergent.
Dominated convergence shows that it is continuous. The same estimate also
makes $\Gamma_\alpha$ continuous as a bilinear map when each domain is given
the graph norm $\|u\|_\infty+\|A_\alpha u\|_\infty$.

For $s>0$, set $f_s=T_s^{(\alpha)}f$ and
$g_s=T_s^{(\alpha)}g$. These functions are smooth, and semigroup theory gives
\[
 f_s\to f,\quad A_\alpha f_s=T_s^{(\alpha)}A_\alpha f\to A_\alpha f
 \quad\text{uniformly},
\]
with the analogous statements for $g$. Applying
\eqref{eq:singint} to the smooth products gives
\[
 A_\alpha(f_sg_s)
 =f_sA_\alpha g_s+g_sA_\alpha f_s-\Gamma_\alpha(f_s,g_s).
\]
The right side converges uniformly, by the preceding graph-norm continuity,
and $f_sg_s\to fg$ uniformly. Since $A_\alpha$ is closed, $fg\in
D(A_\alpha)$ and \eqref{eq:product} holds. This proves
Theorem~\ref{thm:main} for $0<\alpha<2$.

For $\alpha=2$, the generator domain on $C(\T)$ is
$D(A_2)=C^2(\T)$: distributionally $A_2f=-f''$, so continuity of $A_2f$
allows two integrations and yields a classical second derivative. This space
is plainly an algebra. The theorem follows.

For Corollary~\ref{cor:source}, the scaling group of the classical compact
group $\T$ is trivial, and the source identifies
$\D_+=A(\T)\cap P(\T)$, where $A(\T)$ is the Wiener algebra and $P(\T)$ the
positive-definite cone. Every absolutely summable Fourier sequence is a
linear combination of four nonnegative sequences, so
$\operatorname{span}(\D_+)=A(\T)$. Theorem~\ref{thm:main} makes the larger
space $D(L_\alpha)=D(A_\alpha)$ an algebra, hence its intersection with
$A(\T)$ is an algebra as well.

\section{Verification and limitations}

\paragraph{Proof audit.}
The proof uses three checkable inputs: the $L^1$ translation modulus of the
periodic Riesz kernel, the positive periodized singular-integral kernel, and
closedness of the semigroup generator. The exponents in
\eqref{eq:majorant} were checked separately on the three ranges divided by
$\alpha=1$. No computational experiment is used as proof evidence.

\paragraph{Human-review focus.}
The recommended review points are the endpoint estimate
$\rho_1(r)=r\log(e/r)$, graph-norm continuity of $\Gamma_\alpha$, and the
identification $\operatorname{span}(\D_+)=A(\T)$ in the source's conventions.

\paragraph{Scope.}
This does not resolve the source's question for an arbitrary convolution
semigroup, even on $\T$, and does not resolve it for general locally compact
quantum groups. The proof succeeds because the stable symbol gives sharp
potential-kernel regularity matching the singularity of its carr\'e du champ.

\section{Novelty check}

A bounded search checked the exact sentence from Remark 2.9, the source title
together with generator-domain and algebra keywords, and fractional
Laplacian domain-algebra terminology. It located broad Sobolev algebra and
fractional Schauder results, including \cite{Stinga}, but no later paper
stating the source identification or Theorem~\ref{thm:main} on $C(\T)$.
The regularity ingredients are standard, and the theorem may be an unstated
consequence of known Schauder theory. Novelty is therefore \emph{provisional},
not established.

\begin{thebibliography}{9}

\bibitem{SV}
A.~Skalski and A.~Viselter,
\emph{Generating functionals for locally compact quantum groups},
International Mathematics Research Notices 2021 (2021), 132--159;
arXiv:1901.07477.

\bibitem{Stinga}
P.~R.~Stinga,
\emph{User's guide to the fractional Laplacian and the method of semigroups},
in \emph{Handbook of Fractional Calculus with Applications}, Vol.~2,
De Gruyter, 2019, 235--266; arXiv:1808.05159.

\end{thebibliography}

\end{document}
